\section{SİSTEM GEREKSİNİMLERİ}
\label{bolum4}

\subsection{Fonksiyonel Gereksinimler}
\label{fonksiyonel-gereksinimler}

Sistemin fonksiyonel gereksinimleri aşağıdaki tabloda özetlenmektedir:

\begin{table}[H]
\centering
\begin{tabular}{|L{3.5cm}|L{3.5cm}|L{3.5cm}|}
\hline
\textbf{ID} & \textbf{Gereksinim} & \textbf{Açıklama} \\
\hline
FR-1 & Kullanıcı Yönetimi & Kullanıcılar sisteme kayıt olabilmeli, giriş yapabilmeli ve token tabanlı kimlik doğrulama gerçekleştirebilmelidir. \\
\hline
FR-2 & Zafiyet Senaryoları & Güvensiz ortamda Mass Assignment, Broken Authentication, IDOR/BOLA ve BFLA zafiyetleri uygulanmış olmalıdır. \\
\hline
FR-3 & CTF Hedefleri & Kullanıcılar endpoint'ler üzerinden saldırı gerçekleştirerek belirlenmiş hedeflere ulaşabilmelidir. \\
\hline
FR-4 & Güvenli Ortam & Secure API ortamında tüm savunma mekanizmaları aktif ve işlevsel olmalıdır. \\
\hline
FR-5 & Karşılaştırmalı Analiz & Kullanıcılar güvensiz ve güvenli ortamları aynı sistem üzerinde karşılaştırabilmelidir. \\
\hline
FR-6 & Rol Yönetimi & Sistem; user, admin gibi farklı rolleri desteklemeli, rol bazlı erişim kontrolü uygulanmalıdır. \\
\hline
\end{tabular}
\end{table}

\subsection{Güvenlik Gereksinimleri}
\label{g-venlik-gereksinimleri}

Bu proje kapsamında geliştirilen API Fortress sistemi, modern REST API güvenlik prensiplerine uygun şekilde tasarlanmıştır. Sistem içerisinde hem saldırı senaryolarının uygulanabilmesi hem de güvenli backend mimarisinin gösterilebilmesi amacıyla çeşitli güvenlik gereksinimleri belirlenmiştir. Bu gereksinimler, güvenli API ortamında uygulanması zorunlu olan temel savunma mekanizmalarını tanımlamaktadır. Belirlenen güvenlik gereksinimleri; kimlik doğrulama, erişim kontrolü, saldırı filtreleme, parola güvenliği ve servis koruma mekanizmalarını kapsamaktadır.

\noindent\textbf{NFR-1:} Güvenli ortamda tüm endpoint’ler JWT doğrulaması gerektirmelidir.

Güvenli API ortamında yer alan tüm korumalı endpoint’lere erişim sağlanabilmesi için kullanıcıların geçerli bir JWT (JSON Web Token) ile doğrulanması gerekmektedir. JWT tabanlı kimlik doğrulama sistemi, modern REST API mimarilerinde yaygın olarak kullanılan stateless authentication yaklaşımını temsil etmektedir. Kullanıcı sisteme başarılı şekilde giriş yaptıktan sonra backend tarafından imzalanmış bir token üretilmekte ve kullanıcıya iletilmektedir. Sonraki API isteklerinde kullanıcı bu token’ı HTTP Authorization header’ı içerisinde göndermektedir.

Sistem, gelen token’ın doğruluğunu, imzasını, süresini ve kullanıcı bilgilerini kontrol ederek erişim izni vermektedir. Geçersiz, süresi dolmuş veya manipüle edilmiş token’lar sistem tarafından reddedilmektedir. Böylece yetkisiz kullanıcıların korumalı API endpoint’lerine erişmesi engellenmektedir. Bu yapı sayesinde modern backend sistemlerinde kullanılan güvenli kimlik doğrulama mantığı uygulamalı olarak gösterilmektedir.

Ayrıca JWT doğrulaması sayesinde backend tarafında oturum bilgisinin sunucu belleğinde tutulmasına gerek kalmamakta, böylece sistem daha ölçeklenebilir ve dağıtık mimarilere uygun hale gelmektedir. Bu güvenlik gereksinimi özellikle kullanıcı verileri, yönetici işlemleri ve hassas endpoint’lerin korunması açısından kritik öneme sahiptir.

\noindent\textbf{NFR-2:} Rate limiting mekanizması, dakika başına istek sayısını sınırlandırmalıdır.

Sistem içerisinde brute force saldırıları, spam istekler ve servis yoğunlaştırma girişimlerini azaltmak amacıyla rate limiting mekanizması uygulanmıştır. Bu mekanizma, belirli bir istemcinin belirli süre içerisinde gönderebileceği maksimum istek sayısını sınırlandırmaktadır.

Örneğin bir kullanıcı kısa süre içerisinde sürekli login isteği gönderirse sistem bu davranışı şüpheli olarak değerlendirmekte ve belirlenen limit aşıldığında ilgili istemcinin yeni istekleri geçici olarak engellenmektedir. Böylece parola tahmin saldırıları ve otomatik saldırı araçlarının etkisi azaltılmaktadır.

Rate limiting sistemi özellikle aşağıdaki durumlarda kritik koruma sağlamaktadır:

\begin{itemize}[leftmargin=*]
    \item Brute force login saldırıları
    \item API spam istekleri
    \item Kaynak tüketimini artırmaya yönelik saldırılar
    \item Otomatik bot aktiviteleri
    \item Servis yoğunlaştırma girişimleri
\end{itemize}

Bu proje kapsamında Flask-Limiter mekanizması kullanılarak istemci IP adresi bazlı istek kontrolü uygulanmıştır. Böylece kullanıcıların belirlenen limitlerin üzerinde istek göndermesi engellenmiştir. Güvenli backend mimarilerinde rate limiting mekanizmaları servis sürekliliği ve sistem kararlılığı açısından önemli güvenlik katmanlarından biridir.

\noindent\textbf{NFR-3:} WAF sistemi SQL Injection ve XSS içeren istekleri engelleyebilmelidir.

Güvenli API ortamında temel seviyede bir WAF (Web Application Firewall) sistemi uygulanmıştır. Bu sistemin amacı, kullanıcıdan gelen zararlı girdileri analiz ederek potansiyel saldırı girişimlerini engellemektir.

Sistem özellikle aşağıdaki saldırı türlerini tespit etmeye odaklanmaktadır:

\begin{itemize}[leftmargin=*]
    \item SQL Injection
    \item Cross-Site Scripting (XSS)
    \item Zararlı script girişimleri
    \item Şüpheli SQL sorgu yapıları
    \item Veri manipülasyonu girişimleri
\end{itemize}

WAF sistemi, gelen HTTP request içeriklerini regex tabanlı filtreleme mekanizmalarıyla incelemektedir. Örneğin:

\begin{lstlisting}
OR 1=1UNION SELECTDROP TABLE<script>
\end{lstlisting}

gibi zararlı ifadeler tespit edildiğinde sistem isteği engellemekte ve kullanıcıya hata cevabı döndürmektedir.

Bu yaklaşım gerçek kurumsal WAF sistemlerinin basitleştirilmiş bir simülasyonu niteliğindedir. Amaç yalnızca saldırıyı engellemek değil; aynı zamanda kullanıcıların modern backend güvenlik filtreleme mantığını uygulamalı olarak anlayabilmesini sağlamaktır.

WAF mekanizması özellikle input validation süreçlerinin önemini göstermektedir. Güvensiz backend sistemlerinde filtrelenmeyen kullanıcı girdileri ciddi güvenlik açıklarına neden olabilirken; güvenli sistemlerde giriş verilerinin kontrol edilmesi saldırı yüzeyini önemli ölçüde azaltmaktadır.

\noindent\textbf{NFR-4:} Kullanıcı parolaları düz metin olarak saklanmamalı; bcrypt ile hashlenerek depolanmalıdır.

Sistem güvenliği açısından kullanıcı parolalarının düz metin (plain text) olarak saklanması kesinlikle yasaktır. Güvenli API ortamında kullanıcı şifreleri bcrypt algoritması kullanılarak hashlenmekte ve yalnızca hash değerleri veritabanında saklanmaktadır.

Bu yaklaşımın temel amacı, veritabanı sızıntısı durumunda kullanıcı şifrelerinin doğrudan ele geçirilmesini engellemektir. Eğer şifreler düz metin olarak saklanırsa veritabanına erişen saldırganlar tüm kullanıcı hesaplarını kolayca ele geçirebilir. Ancak bcrypt ile hashlenmiş parolalar geri döndürülemez yapıda olduğu için saldırganlar gerçek şifreleri doğrudan elde edememektedir.

\subsubsection{Bcrypt algoritması aynı zamanda:}

\begin{itemize}[leftmargin=*]
    \item Salt mekanizması kullanır,
    \item Rainbow table saldırılarını zorlaştırır,
    \item Brute force saldırı maliyetini artırır,
    \item Her parola için farklı hash üretir.
\end{itemize}

Kullanıcı giriş işlemi sırasında sistem, girilen şifreyi yeniden hashleyerek veritabanındaki hash değeriyle karşılaştırmaktadır. Böylece gerçek parola hiçbir zaman doğrudan saklanmamakta veya sistem içerisinde açık biçimde tutulmamaktadır.

Bu güvenlik gereksinimi modern kimlik doğrulama sistemlerinin temel prensiplerinden biridir ve güvenli backend geliştirme süreçlerinde kritik öneme sahiptir.

\noindent\textbf{NFR-5:} Rol tabanlı yetkilendirme, kullanıcı ve yönetici işlevlerini birbirinden kesin biçimde ayırmalıdır.

Sistem içerisinde Role-Based Access Control (RBAC) yaklaşımı uygulanmıştır. Bu yapı sayesinde kullanıcıların sistem içerisindeki yetkileri roller üzerinden kontrol edilmektedir.

\subsubsection{Projede temel olarak iki farklı rol bulunmaktadır:}

\begin{itemize}[leftmargin=*]
    \item Normal kullanıcı (user)
    \item Yönetici (admin)
\end{itemize}

Her rol yalnızca kendisine izin verilen işlemleri gerçekleştirebilmektedir. Örneğin normal kullanıcılar yalnızca kendi verilerine erişebilirken; yönetici kullanıcılar sistem yönetimi, kullanıcı listeleme veya kritik işlemleri gerçekleştirebilmektedir.

\subsubsection{RBAC sisteminin temel amacı:}

\begin{itemize}[leftmargin=*]
    \item Yetkisiz erişimleri engellemek,
    \item Ayrıcalık kontrolü sağlamak,
    \item Hassas işlemleri korumak,
    \item Kullanıcı izolasyonu oluşturmaktır.
\end{itemize}

Sistem, her API isteğinde kullanıcının token bilgilerini analiz ederek ilgili rolü kontrol etmektedir. Eğer kullanıcı gerekli yetkiye sahip değilse sistem erişimi reddetmektedir. Böylece kullanıcıların yönetici işlemlerine erişmesi veya başka kullanıcıların verilerini görüntülemesi engellenmektedir.

Bu güvenlik gereksinimi özellikle Broken Access Control ve Broken Object Level Authorization (BOLA) gibi modern API zafiyetlerini önlemek açısından büyük önem taşımaktadır. Çünkü birçok gerçek dünya saldırısı yetersiz erişim kontrolü mekanizmalarından kaynaklanmaktadır.

Sonuç olarak bu güvenlik gereksinimleri, API Fortress sisteminin modern REST API güvenlik prensiplerine uygun biçimde çalışmasını sağlamak amacıyla belirlenmiştir. JWT doğrulama, rate limiting, WAF filtreleme, bcrypt tabanlı parola güvenliği ve rol tabanlı erişim kontrolü gibi mekanizmalar sayesinde sistem yalnızca saldırı senaryolarını göstermekle kalmamakta; aynı zamanda güvenli backend geliştirme yaklaşımını da uygulamalı olarak sergilemektedir.

\subsection{Teknik Gereksinimler}
\label{teknik-gereksinimler}

API Fortress sisteminin geliştirme sürecinde modern backend mimarisi, sürdürülebilir yazılım geliştirme yaklaşımı ve güvenli sistem tasarımı dikkate alınarak çeşitli teknik gereksinimler belirlenmiştir. Bu gereksinimler; sistemin geliştirme altyapısını, backend mimarisini, veritabanı yönetimini, konteyner yapısını ve güvenli yapılandırma süreçlerini tanımlamaktadır. Belirlenen teknik gereksinimler aynı zamanda projenin modüler, taşınabilir, sürdürülebilir ve gerçek dünya backend sistemlerine yakın bir yapıda geliştirilmesini amaçlamaktadır.

\noindent\textbf{TR-1:} Backend geliştirme için Python 3.9+ ve Flask framework kullanılmalıdır.

Sistemin backend katmanı Python programlama dili kullanılarak geliştirilmiştir. Python’un tercih edilmesinin temel nedenleri arasında hızlı geliştirme imkânı sunması, geniş kütüphane desteğine sahip olması, okunabilir sözdizimi ve modern web teknolojileriyle uyumlu yapısı yer almaktadır. Özellikle siber güvenlik, veri işleme ve backend geliştirme alanlarında Python yaygın olarak kullanılmaktadır.

Projede minimum Python 3.9 sürümü hedeflenmiştir. Bunun nedeni modern Python sürümleriyle birlikte gelen performans iyileştirmeleri, gelişmiş type hint desteği, güvenlik güncellemeleri ve güncel kütüphanelerle uyumluluktur. Ayrıca Docker ortamında çalışan backend servislerinin stabil biçimde yönetilebilmesi açısından güncel Python sürümleri tercih edilmiştir.

Backend framework olarak Flask kullanılmıştır. Flask, hafif (lightweight) ve esnek yapısı sayesinde REST API geliştirme süreçleri için oldukça uygun bir framework’tür. Flask’ın tercih edilme nedenleri şunlardır:

\begin{itemize}[leftmargin=*]
    \item Minimal ve modüler yapı sunması
    \item REST API geliştirmeye uygun olması
    \item Blueprint mimarisini desteklemesi
    \item Kolay genişletilebilir olması
    \item Güvenlik ve middleware entegrasyonlarının kolay yapılabilmesi
    \item Eğitim ve laboratuvar ortamları için anlaşılır yapı sağlaması
\end{itemize}

Ayrıca Flask ekosistemi sayesinde JWT doğrulama, rate limiting, ORM yönetimi ve CORS entegrasyonu gibi güvenlik bileşenleri sisteme kolayca entegre edilebilmiştir. Böylece proje hem modern backend geliştirme yaklaşımını hem de güvenlik odaklı API mimarisini uygulamalı olarak gösterebilecek bir altyapıya kavuşmuştur.

\noindent\textbf{TR-2:} Veritabanı işlemleri SQLAlchemy ORM üzerinden yürütülmelidir.

Sistem içerisinde veritabanı işlemleri doğrudan ham SQL sorguları ile değil, SQLAlchemy ORM (Object Relational Mapping) katmanı üzerinden yürütülmektedir. ORM yaklaşımının temel amacı veritabanı işlemlerini nesne yönelimli programlama mantığıyla yönetebilmektir.

SQLAlchemy sayesinde veritabanındaki tablolar Python sınıflarıyla temsil edilmektedir. Örneğin kullanıcı tablosu bir User modeli üzerinden tanımlanmakta ve kayıt işlemleri Python nesneleri üzerinden gerçekleştirilmektedir. Böylece geliştiriciler doğrudan karmaşık SQL sorguları yazmak yerine daha okunabilir ve sürdürülebilir kod yapıları oluşturabilmektedir.

\subsubsection{ORM kullanımının sağladığı temel avantajlar şunlardır:}

\begin{itemize}[leftmargin=*]
    \item Kod okunabilirliğinin artması
    \item Veritabanı bağımsızlığı sağlanması
    \item SQL sorgularının merkezi yönetilebilmesi
    \item Model tabanlı veri yönetimi
    \item Daha güvenli sorgu yapıları oluşturulması
    \item Geliştirme sürecinin hızlanması
\end{itemize}

Ayrıca SQLAlchemy parameterized query yaklaşımı kullandığı için SQL Injection riskinin azaltılmasına katkı sağlamaktadır. Güvensiz API ortamında bazı endpoint’lerde bilinçli olarak zafiyet oluşturulsa da güvenli backend ortamında ORM yapısı güvenli veri işleme mantığını temsil etmektedir.

Projede SQLAlchemy kullanılması aynı zamanda gerçek dünya backend mimarilerine daha yakın bir yapı oluşturulmasını sağlamaktadır. Günümüzde birçok profesyonel Python backend sistemi ORM tabanlı veri yönetimi kullanmaktadır.

\noindent\textbf{TR-3:} Sistem Docker konteyner ortamında çalıştırılabilir olmalıdır.

API Fortress sistemi Docker konteyner mimarisi kullanılarak geliştirilmiştir. Docker teknolojisinin kullanılmasının temel amacı sistemin taşınabilir, izole ve tekrar üretilebilir bir laboratuvar ortamı sunmasını sağlamaktır.

Geleneksel geliştirme süreçlerinde farklı işletim sistemleri, bağımlılık sürümleri ve ortam yapılandırmaları çeşitli uyumsuzluk problemlerine neden olabilmektedir. Docker konteyner teknolojisi sayesinde uygulamanın ihtiyaç duyduğu tüm bağımlılıklar aynı konteyner içerisinde paketlenmekte ve sistemin farklı ortamlarda aynı şekilde çalışması sağlanmaktadır.

Docker kullanımının projeye sağladığı avantajlar şunlardır:

\begin{itemize}[leftmargin=*]
    \item Platform bağımsız çalışabilme
    \item Hızlı kurulum ve dağıtım
    \item İzole çalışma ortamı
    \item Tekrar üretilebilir laboratuvar yapısı
    \item Güvensiz servislerin kontrollü şekilde çalıştırılması
    \item Geliştirme ortamı standartlaşması
    \item Mikro servis mimarisine uygun yapı
\end{itemize}

Projede hem güvenli hem de güvensiz backend servisleri ayrı konteynerler içerisinde çalıştırılabilmektedir. Böylece kullanıcılar saldırı testlerini kontrollü laboratuvar ortamında gerçekleştirebilmekte ve sistemin gerçek cihazlara zarar verme riski azaltılmaktadır.

Docker kullanımı aynı zamanda eğitim erişilebilirliğini de artırmaktadır. Kullanıcılar sistemi yalnızca birkaç komut ile kendi bilgisayarlarında çalıştırabilmekte ve karmaşık kurulum süreçleriyle uğraşmadan laboratuvar ortamına erişebilmektedir.

\noindent\textbf{TR-4:} API endpoint’leri Flask Blueprint mimarisiyle modüler olarak organize edilmelidir.

Sistemin backend mimarisinde Flask Blueprint yapısı kullanılmıştır. Blueprint mimarisinin temel amacı büyük ölçekli backend sistemlerini modüler hale getirmek ve farklı işlevleri birbirinden ayrılmış dosyalar halinde organize edebilmektir.

Projede kullanıcı işlemleri, kimlik doğrulama işlemleri ve veri yönetimi gibi farklı API servisleri ayrı route dosyalarına bölünmüştür. Örneğin:

\begin{lstlisting}
routes_auth.pyroutes_users.pyroutes_items.py
\end{lstlisting}

gibi modüller birbirinden bağımsız şekilde yapılandırılmıştır.

Blueprint mimarisinin kullanılmasının temel avantajları şunlardır:

\begin{itemize}[leftmargin=*]
    \item Kod karmaşıklığını azaltmak
    \item Büyük projelerde okunabilirliği artırmak
    \item Endpoint yönetimini kolaylaştırmak
    \item Modüller arası bağımsızlık sağlamak
    \item Yeni servislerin kolay eklenebilmesi
    \item Test süreçlerini kolaylaştırmak
    \item Bakım süreçlerini hızlandırmak
\end{itemize}

Bu yapı sayesinde backend sistemi tek dosyada çalışan karmaşık bir yapı olmaktan çıkarılarak profesyonel backend mimarilerine benzer şekilde organize edilmiştir. Ayrıca modüler yapı sayesinde güvenli ve güvensiz backend ortamları daha kolay yönetilebilmektedir.

Blueprint mimarisi aynı zamanda REST API geliştirme süreçlerinde ölçeklenebilirlik açısından önemli avantaj sağlamaktadır. Yeni endpoint’lerin veya güvenlik mekanizmalarının sisteme eklenmesi mevcut yapıyı bozmadan gerçekleştirilebilmektedir.

\noindent\textbf{TR-5:} Güvenli ortamda gizli bilgiler (secret key, DB şifresi) environment variable olarak yönetilmelidir.

Sistem güvenliği açısından kritik yapılandırma bilgilerinin doğrudan kaynak kod içerisine yazılması büyük güvenlik riski oluşturmaktadır. Bu nedenle güvenli backend ortamında gizli bilgiler environment variable mekanizması kullanılarak yönetilmektedir.

Environment variable yapısı sayesinde aşağıdaki bilgiler kod içerisine gömülmeden yönetilebilmektedir:

\begin{itemize}[leftmargin=*]
    \item JWT secret key
    \item Veritabanı kullanıcı adı ve şifresi
    \item Flask secret key
    \item Debug ayarları
    \item API güvenlik anahtarları
    \item Konteyner yapılandırma bilgileri
\end{itemize}

\subsubsection{Bu yaklaşımın temel avantajları şunlardır:}

\begin{itemize}[leftmargin=*]
    \item Hassas bilgilerin kaynak koddan ayrılması
    \item GitHub gibi platformlarda secret sızıntılarının önlenmesi
    \item Farklı ortamlar için ayrı yapılandırma sağlanması
    \item Production ve development ortamlarının ayrıştırılması
    \item Güvenli dağıtım süreçlerinin desteklenmesi
\end{itemize}

Projede .env dosyaları kullanılarak ortam değişkenleri yönetilmiştir. Flask uygulaması başlatılırken bu değişkenler yüklenmekte ve backend servisleri gerekli yapılandırmaları güvenli şekilde kullanabilmektedir.

Bu gereksinim modern backend geliştirme süreçlerinde kritik güvenlik standartlarından biridir. Özellikle gerçek dünya sistemlerinde gizli bilgilerin kaynak kod içerisine yazılması ciddi güvenlik problemlerine neden olabilmektedir. API Fortress içerisinde bu yaklaşım uygulanarak güvenli yapılandırma yönetimi prensibi uygulamalı biçimde gösterilmiştir.

Sonuç olarak belirlenen teknik gereksinimler; sistemin modern backend geliştirme prensiplerine uygun, modüler, güvenli, taşınabilir ve sürdürülebilir bir yapıda geliştirilmesini sağlamaktadır. Python ve Flask tabanlı backend mimarisi, SQLAlchemy ORM kullanımı, Docker konteyner yapısı, Blueprint tabanlı modüler organizasyon ve environment variable yönetimi sayesinde API Fortress hem eğitimsel hem de teknik açıdan gerçek dünya backend sistemlerine yakın bir laboratuvar ortamı sunmaktadır.


\clearpage
